\scanpage{3}
% Lines 1–2. The two left entries are one key.
\[
\begin{aligned}
 A&\quad\text{V.A.}\\
 B&\quad\text{duale}
\end{aligned}
\qquad
 H^*(A)=\underline{\underline{H}}^*(A,\underline{\underline{\Omega}}_A^*)
 \quad\text{\unsure{Coh.} de Hodge-DeRham}
\]
% Lines 3–4
\[
\begin{aligned}
 t_B&=H^1(A,\underline{O}_A)\\
 t_A&=H^0(A,\Omega_A^1)^{\deleted{\unsure{*}}\vee}\simeq H^n(A,\Omega_A^{n-1})
\end{aligned}
\]

% Lines 5–6
\margin{$\deleted{\unsure{\xi|S}}\unsure{S_A}$}
\diagram{\begin{tikzcd}[ampersand replacement=\&,column sep=large]
 0 \arrow[r,phantom,"{\unsure{\relbar|\to}}" description] \& t_A^\vee \arrow[r] \&
 \deleted{\unsure{H}}\unsure{\mathbb{H}|H}^1(A) \arrow[r] \& t_B \arrow[r] \& 0
\end{tikzcd}}
extension canonique (du moins en car. 0 !)

% Lines 7–9
Autre interprétation : Considérons l'extension \uline{universelle} de $A$ par un groupe vectoriel \inserted{$V$}, \deleted{elle} \inserted{une telle extension \deleted{\unsure{se|est}}} se \deleted{\unsure{fait}} \inserted{classifie} par $H^1(A,\deleted{\unsure{\underline{O}_A|\mathcal{O}_A}}V)=\deleted{V}\inserted{\operatorname{Hom}(H^1(A,O_A)^\vee\unsure{\relbar|,|\to}V)}$, donc l'extension universelle est
% Line 10
\diagram{\begin{tikzcd}[ampersand replacement=\&,column sep=large]
 0 \arrow[r,dash] \& \unsure{\mathbb{V}|W}(\deleted{\unsure{H}}t_B^\vee) \arrow[r,-{Stealth[length=2pt,width=2pt]}] \& E \arrow[r,dash] \& A \arrow[r,-{Stealth[length=2pt,width=2pt]}] \& 0
\end{tikzcd}}
et passant aux espaces tangents on trouve une extension
% Line 13
\diagram{\begin{tikzcd}[ampersand replacement=\&,column sep=large]
 0 \arrow[r] \& t_B^\vee \arrow[r] \& t_E \arrow[r] \& t_A \arrow[r] \& 0
\end{tikzcd}}

% Lines 14–16
Je dis que cette extension n'est autre que $S_B$, (ce qui donne en passant
\[
 t_E=H^1(B) \tag*{)}
\]
% Lines 17–19
et que de plus elle est \uline{transposée} de $S_A$, ce qui donne en passant
\[
 \boxed{H^1(A)\ \text{et}\ H^1(B)\ \text{sont duels l'un de l'autre}}
\]
% Line 20
\margin{$\Vert$}Plus \unsure{gén.,|gén,|fin,} il faut \unsure{présumer|prouver} que $H^*(A)$ et $H^*(B)$
% archive sheet number: 2
